% This is text associated with the open textbook "Open Electromagnetics"
% See immediately below for author, date
% <License info goes here (boilerplate, not read by project compiler)>

%ORB TITLE Units     
%ORB AUTHOR 0       # S.W. Ellingson, Virginia Tech (ellingson@vt.edu)
%ORB DATE   20191113
%ORB ABSTRACT SI system of units, prefixes (i.e., "kilo-", "milli-" etc), ref to 

\label{m0072_Units} % <-- Must match name of this file and module directory name.  Don't mess with this!
\index{units|(}

The term ``unit'' refers to the measure used to express a physical quantity.
For example, the mean radius of the Earth is about 6,371,000~meters; in this case, the unit is the meter.

A number like ``6,371,000'' becomes a bit cumbersome to write, so it is common to use a prefix to modify the unit.
For example, the radius of the Earth is more commonly said to be 6371~kilometers, where one kilometer is understood to mean 1000~meters. 
It is common practice to use prefixes, such as ``kilo-,'' that yield values in the range of $0.001$ to $10,000$.  
A list of standard prefixes appears in Table~\ref{m0072_tPrefix}.
%
\begin{table}[b!]
\begin{center}
\begin{tabular}{lll}
Prefix & Abbreviation & Multiply by: \\
\hline
exa & E & $10^{18}$ \\
peta & P & $10^{15}$ \\
tera & T & $10^{12}$ \\
giga & G & $10^{9}$ \\
mega & M & $10^{6}$ \\
kilo & k & $10^{3}$ \\
milli & m & $10^{-3}$ \\
micro & $\mu$ & $10^{-6}$ \\
nano & n & $10^{-9}$ \\
pico & p & $10^{-12}$ \\
femto & f & $10^{-15}$ \\
atto & a & $10^{-18}$ \\
\end{tabular}
\end{center}
\caption{Prefixes used to modify units.}
\label{m0072_tPrefix}
\end{table}

Writing out the names of units can also become tedious.
For this reason, it is common to use standard abbreviations; e.g., ``6731~km'' as opposed to ``6371~kilometers,'' where ``k'' is the standard abbreviation for the prefix ``kilo'' and ``m'' is the standard abbreviation for ``meter.''
A list of commonly-used base units and their abbreviations are shown in Table~\ref{m0072_tUnits}.
%
\begin{table}
\begin{center}
\begin{tabular}{lll}
Unit & Abbreviation & Quantifies: \\
\hline
ampere & A & electric current \\
coulomb & C & electric charge\\
farad & F & capacitance \\
henry & H & inductance \\
hertz & Hz & frequency\\
joule & J & energy \\
meter & m & distance \\
newton & N & force \\
ohm & $\Omega$ & resistance\\
second & s & time\\
tesla & T & magnetic flux density \\
volt & V & electric potential\\
watt & W & power\\
weber & Wb & magnetic flux \\
\end{tabular}
\end{center}
\caption{Some units that are commonly used in electromagnetics.  
}
\label{m0072_tUnits}
\end{table}

To avoid ambiguity, it is important to always indicate the units of a quantity; e.g., writing ``6371~km'' as opposed to ``6371.''
Failure to do so is a common source of error and misunderstandings.  
An example is the expression: 
$$l = 3t$$ 
where $l$ is length and $t$ is time.  It could be that $l$ is in meters and $t$ is in seconds, in which case ``$3$'' really means ``3~m/s.''
However, if it is intended that $l$ is in kilometers and $t$ is in hours, then ``$3$'' really means ``3~km/h,'' and the equation is literally different.
To patch this up, one might write ``$l = 3t$~m/s''; however, note that this does not resolve the ambiguity we just identified -- i.e., we still don't know the units of the constant ``3.''
Alternatively, one might write ``$l=3t$ where $l$ is in meters and $t$ is in seconds,'' which is unambiguous but becomes quite awkward for more complicated expressions.
A better solution is to write instead:
$$l = \left(3~\mbox{m/s}\right)t$$
or even better:
$$l = at ~~~\mbox{where $a=3$~m/s}$$
since this separates the issue of units from the perhaps more-important fact that $l$ is proportional to $t$ and the constant of proportionality ($a$) is known.

The meter is the fundamental unit of length in the International System of Units, known by its French acronym ``SI'' 
\index{SI (system of units)}
and sometimes informally referred to as the ``metric system.''  
\index{metric system}

%ORB HIGHLIGHT
\begin{mdframed}[backgroundcolor=yellow!20,nobreak=true] 
In this work, we will use SI units exclusively.
\end{mdframed}

Although SI is probably the most popular for engineering use overall, other systems remain in common use.  
For example, the English system, 
\index{English system of units}
where the radius of the Earth might alternatively be said to be about 3959~miles, continues to be used in various applications and to a lesser or greater extent in various regions of the world.
An alternative system in common use in physics and material science applications is the CGS
\index{CGS (system of units)}
(``centimeter-gram-second'') system. The CGS system is similar to SI, but with some significant differences. For example, the base unit of energy in the CGS system is not the ``joule'' but rather the ``erg,'' and the values of some physical constants become unitless. Therefore -- once again -- it is very important to include units whenever values are stated.

SI defines seven fundamental units from which all other units can be derived.  These fundamental units are
distance in meters (m),
time in seconds (s), 
current in amperes (A),
mass in kilograms (kg),
temperature in kelvin (K),
particle count in moles (mol), and
luminosity in candela (cd).
SI units for electromagnetic quantities such as coulombs (C) for charge and volts (V) for electric potential are derived from these fundamental units.

A frequently-overlooked feature of units is their ability to assist in error-checking mathematical expressions.
For example, the electric field intensity may be specified in volts per meter (V/m), so an expression for the electric field intensity that yields units of V/m is said to be ``dimensionally correct'' (but not necessarily correct), whereas an expression that cannot be reduced to units of V/m \emph{cannot} be correct. 

%%%%%%%%%%%%%%%%%%%%%%%%%%%%%%%%%%%%

{\bf Additional Reading:}
\begin{itemize}
\item ``\href{https://en.wikipedia.org/wiki/International_System_of_Units}{International System of Units}'' on Wikipedia.
\index{SI (system of units)}
\item ``\href{https://en.wikipedia.org/wiki/Centimetre-gram-second_system_of_units}{Centimetre-gram-second system of units}'' on Wikipedia.
\index{CGS (system of units)}
\end{itemize}

\index{units|)}
